\documentclass[a4paper,11pt]{report}

\usepackage{amsmath}
\usepackage{graphicx}
\usepackage{biblatex}
\usepackage{listings}
\usepackage{tabularray}
\usepackage{xcolor}
\usepackage{xurl}
\usepackage{hyperref}

% Set table of contents depth
\setcounter{tocdepth}{1}

% Configure references
\addbibresource{../card_usersmanual.bib}
\graphicspath{ {../images/} }

\DefTblrTemplate{contfoot-text}{default}{Continued on the next page}
\DefTblrTemplate{conthead-text}{default}{(Continued)}

% Configure listings
\lstset{
  basicstyle=\ttfamily,
  columns=fullflexible,
  frame=single,
  breaklines=true,
  postbreak=\mbox{\textcolor{red}{$\hookrightarrow$}\space}
}
\renewcommand{\lstlistingname}{Plain text}

% Description
\title{FeliCa Card User's Manual (Unofficial Edition)}
\author{KIRISHIKI Yudai\thanks{Unknown Technologies Inc.}}
\date{July 24, 2026}

\begin{document}

\maketitle

\section*{Disclaimer}

Every effort has been made to ensure the accuracy of this document. However, please note that, by the nature of reverse engineering, it may contain inaccurate results.

Unless otherwise noted, this document assumes the following.

\begin{itemize}
    \item Bit order is such that the least significant bit is bit 0.
    \item Data lengths are given in octets.
\end{itemize}

\tableofcontents

\chapter{Introduction}

FeliCa is used for many stored-fare (SF) electronic money schemes and identification cards, most notably transit IC cards. Yet the most important part of its specification, the part concerning cryptography, is kept confidential and is not widely known. The security of a cryptographic system is only assured once it has been made public and scrutinised by a broad range of experts. For this reason, this document reveals the hidden specifications of FeliCa. It draws heavily on JIS X 6319-4\cite{jis-x-6319-4-2016-en}, the FeliCa Card User's Manual Excerpted Edition (hereafter U-MAN)\cite{felica-usersmanual-en}, felica-tool\cite{felica-tool} and Proxmark3\cite{proxmark3}. Readers who want more detail are encouraged to consult these as well.

\chapter{What Is FeliCa}

FeliCa is a contactless IC card technology developed by Sony Corporation, also known as NFC Type-F. The implementations available on the market are FeliCa Standard and FeliCa Lite-S. This document covers the former.

\chapter{Communication Protocol}

This chapter describes the communication protocol between a card and a Reader/Writer. The physical layer and the data link layer are omitted because they are fully public; only the application layer is covered.

\section{Command Packet}

The data structure of a command packet is shown below. (Table~\ref{table:felica-command-packet-structure})

\begin{tblr}[
        long,
        caption={Data structure of a command packet},
        label={table:felica-command-packet-structure},
    ]{
        colspec={|r|r|l|},
        rowhead=1,
        row{1}={c,m},
        vlines,
        hlines,
    }
    Offset & Length   & Item         \\
    0x00   & 0x01     & Command Code \\
    0x01   & Variable & Command Data
\end{tblr}

\subsection{Command Code}

A one-byte value that identifies the type of the command.

\subsection{Command Data}

Data that specifies how the command is to be processed. Its format differs from command to command.

\section{Response Packet}

The data structure of a response packet is shown below. (Table~\ref{table:felica-response-packet-structure})

\begin{tblr}[
        long,
        caption={Data structure of a response packet},
        label={table:felica-response-packet-structure},
    ]{
        colspec={|r|r|l|},
        rowhead=1,
        row{1}={c,m},
        vlines,
        hlines,
    }
    Offset & Length   & Item          \\
    0x00   & 0x01     & Response Code \\
    0x01   & Variable & Response Data
\end{tblr}

\subsection{Response Code}

An eight-bit value that identifies the type of the response.

\subsection{Response Data}

Data that specifies the result of processing the command. Its format differs from command to command.

\section{Command List}

An overview of each command, together with its Command Code and Response Code, is shown below. (Table~\ref{table:felica-command-list}) In the table, the Command Code is written as CC and the Response Code as RC. Note that commands which are not yet known may exist.

\begin{tblr}[
        long,
        caption={Command list},
        label={table:felica-command-list},
    ]{
        colspec={|l|r|r|X|},
        rowhead=1,
        row{1}={c,m},
        vlines,
        hlines,
    }
    Name                 & CC   & RC   & Overview                                                                                \\
    Polling              & 0x00 & 0x01 & The Reader/Writer captures and identifies a card                                        \\
    Request Service      & 0x02 & 0x03 & Check the existence of Areas and Services and obtain their Key Versions                 \\
    Request Response     & 0x04 & 0x05 & Check the presence and the Mode of a card                                               \\
    {Read Without                                                                                                                \\ Encryption} & 0x06 & 0x07 & Read Block Data from a Service whose Service Attribute requires no authentication \\
    {Write Without                                                                                                               \\ Encryption} & 0x08 & 0x09 & Write Block Data to a Service whose Service Attribute requires no authentication \\
    Search Service Code  & 0x0A & 0x0B & Obtain Area Codes and Service Codes                                                     \\
    Request System Code  & 0x0C & 0x0D & Obtain the System Codes registered in a card                                            \\
    {Request Block                                                                                                               \\ Information} & 0x0E & 0x0F & Obtain the number of Blocks allocated to the specified Nodes \\
    Authentication1      & 0x10 & 0x11 & The Reader/Writer authenticates the card using DES                                      \\
    Authentication2      & 0x12 & 0x13 & The card authenticates the Reader/Writer using DES                                      \\
    Read                 & 0x14 & 0x15 & Read Block Data with DES from a Service whose Service Attribute requires authentication \\
    Write                & 0x16 & 0x17 & Write Block Data with DES to a Service whose Service Attribute requires authentication  \\
    Request Code List    & 0x1A & 0x1B & Iteratively obtain the list of Nodes under the specified parent Node                    \\
    {Request Block                                                                                                               \\ Information Ex} & 0x1E & 0x1F & Obtain the number of Blocks allocated to the specified Nodes and the number of free Blocks \\
    Set Parameter        & 0x20 & 0x21 & Set card communication parameters (encryption scheme and Node Code size)                \\
    {Get Container                                                                                                               \\ Issue Information} & 0x22 & 0x23 & Obtain container information such as the format version and the mobile phone model \\
    {Get Area                                                                                                                    \\ Information} & 0x24 & 0x25 & Obtain information about the specified Area \\
    Get Node Property    & 0x28 & 0x29 & Obtain Node Property                                                                    \\
    {Get Container                                                                                                               \\ Property} & 0x2E & 0x2F & Obtain container properties \\
    Request Service v2   & 0x32 & 0x33 & Check the existence of Areas and Services and obtain their Key Versions (supports AES)  \\
    {Internal Authenticate                                                                                                       \\ and Read} & 0x34 & 0x35 & Perform internal authentication against a Communication-with-MAC-enabled Service and read Block Data \\
    {External Authenticate                                                                                                       \\ and Write} & 0x36 & 0x37 & Perform external authentication against a Communication-with-MAC-enabled Service and write Block Data \\
    Get System Status    & 0x38 & 0x39 & Obtain the configuration state of each System                                           \\
    {Request Product                                                                                                             \\ Information} & 0x3A & 0x3B & Obtain product information \\
    {Request                                                                                                                     \\ Specification \\ Version} & 0x3C & 0x3D & Obtain the version of the card OS \\
    Reset Mode           & 0x3E & 0x3F & Reset the Mode to Mode0                                                                 \\
    Authentication1 v2   & 0x40 & 0x41 & The Reader/Writer authenticates the card using AES                                      \\
    Authentication2 v2   & 0x42 & 0x43 & The card authenticates the Reader/Writer using AES                                      \\
    Read v2              & 0x44 & 0x45 & Read Block Data with AES from a Service whose Service Attribute requires authentication \\
    Write v2             & 0x46 & 0x47 & Write Block Data with AES to a Service whose Service Attribute requires authentication  \\
    Delete Key           & 0x48 & 0x49 & Stop the use of the DES key on a Node that has both an AES key and a DES key            \\
    Update Random ID     & 0x4C & 0x4D & Update the random ID (IDr)                                                              \\
    Get Container ID     & 0x70 & 0x71 & Obtain the container ID from a mobile FeliCa device                                     \\
    Set Node Property    & 0x78 & 0x79 & Set Node Property                                                                       \\
    Register Issue ID    & 0x80 & 0x81 & Initialise a System and register its issue ID                                           \\
    Register Area        & 0x82 & 0x83 & Register an Area                                                                        \\
    Register Service     & 0x84 & 0x85 & Register a Service                                                                      \\
    Register Issue ID Ex & 0x86 & 0x87 & Initialise a System and register its issue ID                                           \\
    Change System Block  & 0x8E & 0x8F & Commit the results of issuance commands
\end{tblr}

\section{Manufacture ID and Manufacture Parameter}

This section describes the manufacture ID (IDm) and the manufacture parameter (PMm), which are obtained as the response data of the Polling command.

\subsection{IDm}

The IDm is the ID with which a Reader/Writer identifies a card. When a card contains several Systems, each System has a different IDm.

The data structure of the IDm is shown below. (Table~\ref{table:felica-idm-structure}) Note that the upper four bits of the first byte of the manufacturer code indicate the System number within the card.

\begin{tblr}[
        long,
        caption={Data structure of the IDm},
        label={table:felica-idm-structure},
    ]{
        colspec={|r|r|l|},
        rowhead=1,
        row{1}={c,m},
        vlines,
        hlines,
    }
    Offset & Length & Item                       \\
    0x00   & 0x02   & Manufacturer code          \\
    0x02   & 0x06   & Card identification number
\end{tblr}

\subsection{PMm}

The data structure of the PMm is shown below. (Table~\ref{table:felica-pmm-structure}) The ROM type and the IC type together are called the IC code.

\begin{tblr}[
        long,
        caption={Data structure of the PMm},
        label={table:felica-pmm-structure},
    ]{
        colspec={|r|r|l|},
        rowhead=1,
        row{1}={c,m},
        vlines,
        hlines,
    }
    Offset & Length & Item                            \\
    0x00   & 0x01   & ROM type                        \\
    0x01   & 0x01   & IC type                         \\
    0x02   & 0x06   & Maximum response time parameter
\end{tblr}

\chapter{File System}

Internally, FeliCa has a hierarchical structure consisting of Systems, Areas, Services and Blocks. A System contains Areas, an Area contains Sub-Areas or Services, and a Service contains Blocks. Any of these may exist more than once on a single card. Of these, Systems, Areas and Services are metadata resembling directories, while Blocks are the regions that hold the actual data. Collectively they are called Nodes. Each Node has a code that represents it. The existence of any Node consumes Blocks.

\section{System}

A System is a logical card unit and sits at the top of the hierarchy. The code that represents a System (as a Node) is 0xFFFF. The following System Definition Information is defined for a System.

\begin{itemize}
    \item System Code
    \item Issue ID information
    \item System Key
    \item System Key Version
    \item System property
\end{itemize}

\subsection{System Code}

A System Code is a two-byte value. Confusingly, the code that represents a System (as a Node) and the System Code are different concepts. Examples of System Codes are shown below. (Table~\ref{table:felica-system-code-examples})

\begin{tblr}[
        long,
        caption={Examples of System Codes},
        label={table:felica-system-code-examples},
    ]{
        colspec={|r|l|},
        rowhead=1,
        row{1}={c,m},
        vlines,
        hlines,
    }
    System Code & Name                           \\
    0x0003      & Japan Railway Cybernetics Area \\
    0x80CD      & Free Area                      \\
    0xFE00      & Common Area                    \\
    0xFE0F      & Management Area
\end{tblr}

\subsection{Issue ID Information}

The Issue ID information consists of the issue ID (IDi) and the Issue Parameter (PMi), each eight bytes long. It can be obtained from the response data once mutual authentication has succeeded with the Authentication2 command or the Authentication2 v2 command. Converting the IDi into a string with a particular algorithm yields the ID number printed at the bottom right of the back of a transit IC card conforming to the Japan Railway Cybernetics standard.

\subsection{System Key}

A System Key is an eight-byte value under DES, and a sixteen-byte value under AES.

\subsection{System Key Version}

A System Key Version is a two-byte value.

\subsection{System Property}

Details unknown.

\subsection{Switching between Systems}

When a card receives a Polling command, or receives a command packet addressed to an IDm other than that of the System currently being operated, the System currently being operated is switched and reset to Mode0. However, when the Authentication1, Authentication1 v2 or Internal Authenticate and Read command succeeds and System switching occurs, the card enters Mode1.

\section{Area}

An Area is contained in a System or in at least one parent Area. The following Area Definition Information is defined for an Area.

\begin{itemize}
    \item Area Code
    \item End Service Code
    \item Number of assigned Blocks
    \item Area Key
    \item Area Key Version
    \item Area property
\end{itemize}

\subsection{Area Code}

An Area Code is a two-byte value. Bits 6 to 15 give the Area Number and bits 0 to 5 give the Area Attribute (Table~\ref{table:felica-area-attributes}). The Area Code also indicates the logical start point of the Area within the card.

\begin{tblr}[
        long,
        caption={Area Attributes},
        label={table:felica-area-attributes},
    ]{
        colspec={|l|r|},
        rowhead=1,
        row{1}={c,m},
        vlines,
        hlines,
    }
    Area Attribute                            & Value    \\
    Area in which Sub-Areas can be created    & 0b000000 \\
    Area in which Sub-Areas cannot be created & 0b000001
\end{tblr}

\subsection{End Service Code}

The logical end point of the Area within the card.

\subsection{Number of Assigned Blocks}

The number of Blocks allocated to the Area.

\subsection{Area Key}

An Area Key is an eight-byte value under DES, and a sixteen-byte value under AES.

\subsection{Area Key Version}

An Area Key Version is a two-byte value.

\subsection{Area Property}

Details unknown.

\subsection{Area 0}

A System always contains Area 0, whose range is 0x0000 to 0xFFFE.

\section{Service}

A Service is contained in at least one Area. The following Service Definition Information is defined for a Service.

\begin{itemize}
    \item Service Code
    \item Number of assigned Blocks
    \item Service Key
    \item Service Key Version
    \item Service property
\end{itemize}

\subsection{Service Code}

A Service Code is a two-byte value. Bits 6 to 15 give the Service Number and bits 0 to 5 give the Service Attribute (Table~\ref{table:felica-service-attributes}). The Service Code also indicates the logical position of the Service within the card.

\begin{tblr}[
        long,
        caption={Service Attributes},
        label={table:felica-service-attributes},
    ]{
        colspec={|l|X|r|},
        rowhead=1,
        row{1}={c,m},
        vlines,
        hlines,
    }
    Service Attribute & Access control                                 & Value    \\
    \SetCell[r=4]{l}Random Service
                      & Read/Write Access: authentication required     & 0b001000 \\
                      & Read/Write Access: authentication not required & 0b001001 \\
                      & Read Only Access: authentication required      & 0b001010 \\
                      & Read Only Access: authentication not required  & 0b001011 \\
    \SetCell[r=4]{l}Cyclic Service
                      & Read/Write Access: authentication required     & 0b001100 \\
                      & Read/Write Access: authentication not required & 0b001101 \\
                      & Read Only Access: authentication required      & 0b001110 \\
                      & Read Only Access: authentication not required  & 0b001111 \\
    \SetCell[r=8]{l}Purse Service
                      & Direct Access: authentication required         & 0b010000 \\
                      & Direct Access: authentication not required     & 0b010001 \\
                      & {Cashback/Decrement Access:                               \\ authentication required} & 0b010010 \\
                      & {Cashback/Decrement Access:                               \\ authentication not required} & 0b010011 \\
                      & Decrement Access: authentication required      & 0b010100 \\
                      & Decrement Access: authentication not required  & 0b010101 \\
                      & Read Only Access: authentication required      & 0b010110 \\
                      & Read Only Access: authentication not required  & 0b010111
\end{tblr}

\subsection{Number of Assigned Blocks}

The number of Blocks allocated to the Service.

\subsection{Service Key}

A Service Key is an eight-byte value under DES, and a sixteen-byte value under AES.

\subsection{Service Key Version}

A Service Key Version is a two-byte value.

\subsection{Service Property}

Set only on products equipped with the Value-Limited Purse Service option or the Communication with MAC option.

\subsection{Overlap Service}

Managing the same group of Blocks with more than one Service Code is called overlapping, and a Service that overlaps a Service having the same Service Number within the same System is called an Overlap Service. The figure showing an example of an Overlap Service, in the "Overlap Service" section of U-MAN\cite{felica-usersmanual-en}, aids understanding.

\section{Block}

A Block is the actual data in the non-volatile memory of a card, divided into sixteen-byte units. Apart from those that hold metadata, a fixed number of Blocks is allocated to each Service, and the Area containing the Service manages the total allocation. The figure showing how Blocks are managed by Area, in the "Area" section of U-MAN\cite{felica-usersmanual-en}, aids understanding.

\chapter{Mode}

A card has states called Modes, from Mode0 to Mode3. The commands that the card can execute are restricted by the Mode. When power is supplied after a power-off state, the card enters Mode0. Mode1 and above are divided into DES, AES and Communication with MAC; executing the Authentication1, Authentication1 v2 or Internal Authenticate and Read command respectively takes the card to Mode1.

Under DES and AES, executing the Authentication2 or Authentication2 v2 command in Mode1 takes the card to Mode2 of the corresponding scheme. In Mode2 the commands that require mutual authentication (Read, Write, Read v2, Write v2 and so on) can be executed. Executing an issuance command takes the card to Mode3 of the corresponding scheme.

Communication with MAC, on the other hand, has neither Mode2 nor Mode3. The command that can be executed in Mode1 is External Authenticate and Write, and executing it returns the card to Mode0.

Once the card has left Mode0 it no longer accepts the Polling command. However, a Polling command that specifies another System in order to perform Switching between Systems can be executed in any Mode. The Reset Mode command also returns the card to Mode0 from any Mode. The current Mode can be checked with the Request Response command. For further detail, see the "Mode" section of U-MAN\cite{felica-usersmanual-en}.

\chapter{Commands}

This chapter describes the parameters given to commands and the details of each command. Because of space constraints, the details of some commands are omitted.

\section{Block List and Block List Element}

A Block List is a set of Block List Elements that identify the Service and the Block Number to be accessed. A Block List Element is either two bytes long (Table~\ref{table:felica-block-list-elem-2b-structure}) or three bytes long (Table~\ref{table:felica-block-list-elem-3b-structure}). Note that the offsets and lengths in these two tables are given in bits.

\begin{tblr}[
        long,
        caption={Data structure of a Block List Element (two bytes)},
        label={table:felica-block-list-elem-2b-structure},
    ]{
        colspec={|r|r|X|},
        rowhead=1,
        row{1}={c,m},
        vlines,
        hlines,
    }
    Offset & Length & Item                       \\
    0      & 1      & Length (0b1)               \\
    1      & 3      & Access Mode                \\
    4      & 4      & Service Code List Order    \\
    8      & 8      & Block Number / Key Version
\end{tblr}

\begin{tblr}[
        long,
        caption={Data structure of a Block List Element (three bytes)},
        label={table:felica-block-list-elem-3b-structure},
    ]{
        colspec={|r|r|X|},
        rowhead=1,
        row{1}={c,m},
        vlines,
        hlines,
    }
    Offset & Length & Item                                       \\
    0      & 1      & Length (0b0)                               \\
    1      & 3      & Access Mode                                \\
    4      & 4      & Service Code List Order                    \\
    8      & 16     & Block Number / Key Version (little endian)
\end{tblr}

\subsection{Access Mode}

Specifies how the Node targeted by the Block List Element is accessed. (Table~\ref{table:felica-ble-access-modes})

\begin{tblr}[
        long,
        caption={Access Modes},
        label={table:felica-ble-access-modes},
    ]{
        colspec={|l|r|},
        rowhead=1,
        row{1}={c,m},
        vlines,
        hlines,
    }
    Access Mode                                   & Value \\
    Other than Cashback Access to a Purse Service & 0b000 \\
    Cashback access to a Purse Service            & 0b001 \\
    Key change                                    & 0b100
\end{tblr}

\section{Polling}

The Polling command is used by a Reader/Writer to capture and identify a card. The IDm and PMm of the System having the specified System Code can be obtained.

\subsection{Data Structure of the Command Packet}

The data structure of the Polling command packet is shown below. (Table~\ref{table:felica-polling-command-packet-structure})

\begin{tblr}[
        long,
        caption={Data structure of the Polling command packet},
        label={table:felica-polling-command-packet-structure},
    ]{
        colspec={|r|r|l|X|},
        rowhead=1,
        row{1}={c,m},
        vlines,
        hlines,
    }
    Offset & Length & Item            & Value or remarks                                                                                                                                                                                                \\
    0x00   & 0x01   & Command Code    & 0x00                                                                                                                                                                                                            \\
    0x01   & 0x02   & System Code (x) & 0x0000 $\leq$ x $\leq$ 0xFFFF (0xFF in either byte is a wildcard). Big endian                                                                                                                                   \\
    0x03   & 0x01   & Request Code    & 0x00: no request, 0x01: System Code request, 0x02: communication performance request                                                                                                                            \\
    0x04   & 0x01   & Time Slot       & Specifies the maximum number of Time Slots that may respond. 0x00, 0x01, 0x03, 0x07, 0x0F (for details see the table of Time Slot specifications in the "Polling" section of U-MAN\cite{felica-usersmanual-en})
\end{tblr}

\subsection{Data Structure of the Response Packet}

The data structure of the Polling response packet is shown below. (Table~\ref{table:felica-polling-response-packet-structure})

\begin{tblr}[
        long,
        caption={Data structure of the Polling response packet},
        label={table:felica-polling-response-packet-structure},
    ]{
        colspec={|r|r|l|X|},
        rowhead=1,
        row{1}={c,m},
        vlines,
        hlines,
    }
    Offset & Length & Item          & Value or remarks                                                                                                                 \\
    0x00   & 0x01   & Response Code & 0x01                                                                                                                             \\
    0x01   & 0x08   & IDm           & -                                                                                                                                \\
    0x09   & 0x08   & PMm           & -                                                                                                                                \\
    0x11   & 0x02   & Request Data  & Returned only when the Request Code in the command is other than 0x00 and the specified Request Code is supported by the product
\end{tblr}

If a Request Code that is not supported is specified, no Request Data is appended to the response packet. Implementations must assume that Request Data may not be returned even when a Request Code is specified.

\section{Request Service}

The Request Service command checks the existence of Areas and Services and obtains their Key Versions.

\subsection{Data Structure of the Command Packet}

The data structure of the Request Service command packet is shown below. (Table~\ref{table:felica-request-service-command-packet-structure})

\begin{tblr}[
        long,
        caption={Data structure of the Request Service command packet},
        label={table:felica-request-service-command-packet-structure},
    ]{
        colspec={|r|r|l|X|},
        rowhead=1,
        row{1}={c,m},
        vlines,
        hlines,
    }
    Offset & Length & Item               & Value or remarks          \\
    0x00   & 0x01   & Command Code       & 0x02                      \\
    0x01   & 0x08   & IDm                & -                         \\
    0x09   & 0x01   & Number of Node (n) & 0x01 $\leq$ n $\leq$ 0x20 \\
    0x0A   & 2n     & Node Code List     & Little endian
\end{tblr}

\subsection{Data Structure of the Response Packet}

The data structure of the Request Service response packet is shown below. (Table~\ref{table:felica-request-service-response-packet-structure})

\begin{tblr}[
        long,
        caption={Data structure of the Request Service response packet},
        label={table:felica-request-service-response-packet-structure},
    ]{
        colspec={|r|r|l|X|},
        rowhead=1,
        row{1}={c,m},
        vlines,
        hlines,
    }
    Offset & Length & Item                  & Value or remarks                                 \\
    0x00   & 0x01   & Response Code         & 0x03                                             \\
    0x01   & 0x08   & IDm                   & -                                                \\
    0x09   & 0x01   & Number of Node (n)    & 0x01 $\leq$ n $\leq$ 0x20                        \\
    0x0A   & 2n     & Node Key Version List & 0xFFFF if the Node does not exist. Little endian
\end{tblr}

\section{Request Response}

The Request Response command checks the presence and the Mode of a card.

\subsection{Data Structure of the Command Packet}

The data structure of the Request Response command packet is shown below. (Table~\ref{table:felica-request-response-command-packet-structure})

\begin{tblr}[
        long,
        caption={Data structure of the Request Response command packet},
        label={table:felica-request-response-command-packet-structure},
    ]{
        colspec={|r|r|l|X|},
        rowhead=1,
        row{1}={c,m},
        vlines,
        hlines,
    }
    Offset & Length & Item         & Value or remarks \\
    0x00   & 0x01   & Command Code & 0x04             \\
    0x01   & 0x08   & IDm          & -
\end{tblr}

\subsection{Data Structure of the Response Packet}

The data structure of the Request Response response packet is shown below. (Table~\ref{table:felica-request-response-response-packet-structure})

\begin{tblr}[
        long,
        caption={Data structure of the Request Response response packet},
        label={table:felica-request-response-response-packet-structure},
    ]{
        colspec={|r|r|l|X|},
        rowhead=1,
        row{1}={c,m},
        vlines,
        hlines,
    }
    Offset & Length & Item          & Value or remarks          \\
    0x00   & 0x01   & Response Code & 0x05                      \\
    0x01   & 0x08   & IDm           & -                         \\
    0x09   & 0x01   & Mode (n)      & 0x00 $\leq$ n $\leq$ 0x03
\end{tblr}

\section{Read Without Encryption}

The Read Without Encryption command reads Block Data from a Service whose Service Attribute requires no authentication.

\subsection{Data Structure of the Command Packet}

The data structure of the Read Without Encryption command packet is shown below. (Table~\ref{table:felica-rwoe-command-packet-structure})

\begin{tblr}[
        long,
        caption={Data structure of the Read Without Encryption command packet},
        label={table:felica-rwoe-command-packet-structure},
    ]{
        colspec={|r|r|l|X|},
        rowhead=1,
        row{1}={c,m},
        vlines,
        hlines,
    }
    Offset & Length                & Item                  & Value or remarks                                                          \\
    0x00   & 0x01                  & Command Code          & 0x06                                                                      \\
    0x01   & 0x08                  & IDm                   & -                                                                         \\
    0x09   & 0x01                  & Number of Service (m) & 0x01 $\leq$ m $\leq$ 0x10                                                 \\
    0x0A   & 2m                    & Service Code List     & Little endian                                                             \\
    -      & 0x01                  & Number of Block (n)   & 0x01 $\leq$ n $\leq$ the maximum number of readable Blocks of the product \\
    -      & 2n $\leq$ N $\leq$ 3n & Block List            & -
\end{tblr}

\subsection{Data Structure of the Response Packet}

The data structure of the Read Without Encryption response packet is shown below. (Table~\ref{table:felica-rwoe-response-packet-structure})

\begin{tblr}[
        long,
        caption={Data structure of the Read Without Encryption response packet},
        label={table:felica-rwoe-response-packet-structure},
    ]{
        colspec={|r|r|l|X|},
        rowhead=1,
        row{1}={c,m},
        vlines,
        hlines,
    }
    Offset & Length & Item                & Value or remarks                                                                         \\
    0x00   & 0x01   & Response Code       & 0x07                                                                                     \\
    0x01   & 0x08   & IDm                 & -                                                                                        \\
    0x09   & 0x01   & Status Flag1        & 0x00: normal                                                                             \\
    0x0A   & 0x01   & Status Flag2        & 0x00: normal                                                                             \\
    0x0B   & 0x01   & Number of Block (n) & 0x01 $\leq$ n $\leq$ the maximum number of simultaneously readable Blocks of the product \\
    0x0C   & 16n    & Block Data          & -
\end{tblr}

\section{Write Without Encryption}

The Write Without Encryption command writes Block Data to a Service whose Service Attribute requires no authentication.

\subsection{Data Structure of the Command Packet}

The data structure of the Write Without Encryption command packet is shown below. (Table~\ref{table:felica-wwoe-command-packet-structure})

\begin{tblr}[
        long,
        caption={Data structure of the Write Without Encryption command packet},
        label={table:felica-wwoe-command-packet-structure},
    ]{
        colspec={|r|r|l|X|},
        rowhead=1,
        row{1}={c,m},
        vlines,
        hlines,
    }
    Offset & Length                & Item                  & Value or remarks                                                                         \\
    0x00   & 0x01                  & Command Code          & 0x08                                                                                     \\
    0x01   & 0x08                  & IDm                   & -                                                                                        \\
    0x09   & 0x01                  & Number of Service (m) & 0x01 $\leq$ m $\leq$ 0x10                                                                \\
    0x0A   & 2m                    & Service Code List     & Little endian                                                                            \\
    -      & 0x01                  & Number of Block (n)   & 0x01 $\leq$ n $\leq$ the maximum number of simultaneously writable Blocks of the product \\
    -      & 2n $\leq$ N $\leq$ 3n & Block List            & -                                                                                        \\
    -      & 16n                   & Block Data            & -
\end{tblr}

\subsection{Data Structure of the Response Packet}

The data structure of the Write Without Encryption response packet is shown below. (Table~\ref{table:felica-wwoe-response-packet-structure})

\begin{tblr}[
        long,
        caption={Data structure of the Write Without Encryption response packet},
        label={table:felica-wwoe-response-packet-structure},
    ]{
        colspec={|r|r|l|X|},
        rowhead=1,
        row{1}={c,m},
        vlines,
        hlines,
    }
    Offset & Length & Item          & Value or remarks \\
    0x00   & 0x01   & Response Code & 0x09             \\
    0x01   & 0x08   & IDm           & -                \\
    0x09   & 0x01   & Status Flag1  & 0x00: normal     \\
    0x0A   & 0x01   & Status Flag2  & 0x00: normal
\end{tblr}

\section{Search Service Code}

The Search Service Code command obtains Area Codes and Service Codes.

\subsection{Data Structure of the Command Packet}

The data structure of the Search Service Code command packet is shown below. (Table~\ref{table:felica-search-service-code-command-packet-structure})

\begin{tblr}[
        long,
        caption={Data structure of the Search Service Code command packet},
        label={table:felica-search-service-code-command-packet-structure},
    ]{
        colspec={|r|r|l|X|},
        rowhead=1,
        row{1}={c,m},
        vlines,
        hlines,
    }
    Offset & Length & Item         & Value or remarks                             \\
    0x00   & 0x01   & Command Code & 0x0A                                         \\
    0x01   & 0x08   & IDm          & -                                            \\
    0x09   & 0x02   & Node Index   & 0x0000 $\leq$ n $\leq$ 0xFFFF. Little endian
\end{tblr}

\subsection{Data Structure of the Response Packet}

The data structure of the Search Service Code response packet is shown below. (Table~\ref{table:felica-search-service-code-response-packet-structure})

\begin{tblr}[
        long,
        caption={Data structure of the Search Service Code response packet},
        label={table:felica-search-service-code-response-packet-structure},
    ]{
        colspec={|r|r|l|X|},
        rowhead=1,
        row{1}={c,m},
        vlines,
        hlines,
    }
    Offset & Length       & Item          & Value or remarks                                                                   \\
    0x00   & 0x01         & Response Code & 0x0B                                                                               \\
    0x01   & 0x08         & IDm           & -                                                                                  \\
    0x09   & 0x02 or 0x04 & Node Code     & {Two bytes: a Service Code. Four bytes: an Area Code and an Area End Service Code. \\ Both are little endian. 0xFFFF indicates that no corresponding Node exists}
\end{tblr}

The Nodes in a card can be enumerated by executing this command while incrementing the Node index from 0, treating the point at which the Node code in the response becomes 0xFFFF as the end.

\section{Request System Code}

The Request System Code command obtains the System Codes registered in a card.

\subsection{Data Structure of the Command Packet}

The data structure of the Request System Code command packet is shown below. (Table~\ref{table:felica-request-system-code-command-packet-structure})

\begin{tblr}[
        long,
        caption={Data structure of the Request System Code command packet},
        label={table:felica-request-system-code-command-packet-structure},
    ]{
        colspec={|r|r|l|X|},
        rowhead=1,
        row{1}={c,m},
        vlines,
        hlines,
    }
    Offset & Length & Item         & Value or remarks \\
    0x00   & 0x01   & Command Code & 0x0C             \\
    0x01   & 0x08   & IDm          & -
\end{tblr}

\subsection{Data Structure of the Response Packet}

The data structure of the Request System Code response packet is shown below. (Table~\ref{table:felica-request-system-code-response-packet-structure})

\begin{tblr}[
        long,
        caption={Data structure of the Request System Code response packet},
        label={table:felica-request-system-code-response-packet-structure},
    ]{
        colspec={|r|r|l|X|},
        rowhead=1,
        row{1}={c,m},
        vlines,
        hlines,
    }
    Offset & Length & Item                      & Value or remarks                              \\
    0x00   & 0x01   & Response Code             & 0x0D                                          \\
    0x01   & 0x08   & IDm                       & -                                             \\
    0x09   & 0x01   & Number of System Code (n) & -                                             \\
    0x0A   & 2n     & System Code List          & Enumerated in order from System 0. Big endian
\end{tblr}

\section{Request Block Information}

The Request Block Information command obtains the number of Blocks allocated to the specified Nodes. It is supported by mobile FeliCa only.

\subsection{Data Structure of the Command Packet}

The data structure of the Request Block Information command packet is shown below. (Table~\ref{table:felica-request-block-information-command-packet-structure})

\begin{tblr}[
        long,
        caption={Data structure of the Request Block Information command packet},
        label={table:felica-request-block-information-command-packet-structure},
    ]{
        colspec={|r|r|l|X|},
        rowhead=1,
        row{1}={c,m},
        vlines,
        hlines,
    }
    Offset & Length & Item                    & Value or remarks \\
    0x00   & 0x01   & Command Code            & 0x0E             \\
    0x01   & 0x08   & IDm                     & -                \\
    0x09   & 0x01   & Number of Node Code (n) & -                \\
    0x0A   & 2n     & Node Code List          & Little endian
\end{tblr}

\subsection{Data Structure of the Response Packet}

The data structure of the Request Block Information response packet is shown below. (Table~\ref{table:felica-request-block-information-response-packet-structure})

\begin{tblr}[
        long,
        caption={Data structure of the Request Block Information response packet},
        label={table:felica-request-block-information-response-packet-structure},
    ]{
        colspec={|r|r|l|X|},
        rowhead=1,
        row{1}={c,m},
        vlines,
        hlines,
    }
    Offset & Length & Item                            & Value or remarks                                  \\
    0x00   & 0x01   & Response Code                   & 0x0F                                              \\
    0x01   & 0x08   & IDm                             & -                                                 \\
    0x09   & 0x01   & Number of Block Information (n) & -                                                 \\
    0x0A   & 2n     & Block Information List          & Gives the number of Blocks allocated to each Node
\end{tblr}

\section{Authentication1}

The Authentication1 command starts the card-side authentication challenge under DES.

\subsection{Data Structure of the Command Packet}

For the data structure of the Authentication1 command packet, see Table~\ref{table:felica-auth1-command-packet-structure}.

\subsection{Data Structure of the Response Packet}

For the data structure of the Authentication1 response packet, see Table~\ref{table:felica-auth1-response-packet-structure}.

\section{Authentication2}

The Authentication2 command completes DES mutual authentication.

\subsection{Data Structure of the Command Packet}

For the data structure of the Authentication2 command packet, see Table~\ref{table:felica-auth2-command-packet-structure}.

\subsection{Data Structure of the Response Packet}

For the data structure of the Authentication2 response packet, see Table~\ref{table:felica-auth2-response-packet-structure}.

\section{Read}

The Read command reads Block Data in a DES session established after mutual authentication.

\subsection{Data Structure of the Command Packet}

For the structure of the internal payload of the Read command before encryption, see Table~\ref{table:felica-secure-read-common-command-payload-structure}.

\subsection{Data Structure of the Response Packet}

For the structure of the internal payload of the Read response before encryption, see Table~\ref{table:felica-secure-read-common-response-payload-structure}.

\section{Write}

The Write command writes Block Data in a DES session established after mutual authentication.

\subsection{Data Structure of the Command Packet}

For the structure of the internal payload of the Write command before encryption, see Table~\ref{table:felica-secure-write-common-command-payload-structure}.
For the key change package stored in the Block Data when the Access Mode is 0b100 (key change), see Table~\ref{table:felica-secure-key-change-package-structure} and the generation algorithm in the same section.

\subsection{Data Structure of the Response Packet}

For the structure of the internal payload of the Write response before encryption, see Table~\ref{table:felica-secure-write-common-response-payload-structure}.

\section{Get Node Property}

The Get Node Property command obtains Node Property. The Node Property of the Value-Limited Purse Service option or of the Communication with MAC option can be obtained. This command is implemented only on some AES card products and AES/DES card products.

\subsection{Data Structure of the Command Packet}

The data structure of the Get Node Property command packet is shown below. (Table~\ref{table:felica-get-node-property-command-packet-structure})

\begin{tblr}[
        long,
        caption={Data structure of the Get Node Property command packet},
        label={table:felica-get-node-property-command-packet-structure},
    ]{
        colspec={|r|r|l|X|},
        rowhead=1,
        row{1}={c,m},
        vlines,
        hlines,
    }
    Offset & Length & Item               & Value or remarks                                                                \\
    0x00   & 0x01   & Command Code       & 0x28                                                                            \\
    0x01   & 0x08   & IDm                & -                                                                               \\
    0x09   & 0x01   & Target             & 0x00: Value-Limited Purse Service, 0x01: Communication-with-MAC-enabled Service \\
    0x0A   & 0x01   & Number of Node (n) & 0x01 $\leq$ n $\leq$ 0x10                                                       \\
    0x0B   & 2n     & Node Code List     & Little endian
\end{tblr}

\subsection{Data Structure of the Response Packet}

The data structure of the Get Node Property response packet is shown below. (Table~\ref{table:felica-get-node-property-response-packet-structure})

\begin{tblr}[
        long,
        caption={Data structure of the Get Node Property response packet},
        label={table:felica-get-node-property-response-packet-structure},
    ]{
        colspec={|r|r|l|X|},
        rowhead=1,
        row{1}={c,m},
        vlines,
        hlines,
    }
    Offset & Length & Item               & Value or remarks                                                                         \\
    0x00   & 0x01   & Response Code      & 0x29                                                                                     \\
    0x01   & 0x08   & IDm                & -                                                                                        \\
    0x09   & 0x01   & Status Flag1       & 0x00: normal                                                                             \\
    0x0A   & 0x01   & Status Flag2       & 0x00: normal                                                                             \\
    0x0B   & 0x01   & Number of Node (n) & Returned only when Status Flag1 is 0x00                                                  \\
    0x0C   & mn     & Node Property      & {Returned only when Status Flag1 is 0x00. Enumerated in the order of the Node Code List. \\ m = 10 when the target is 0x00, m = 1 when the target is 0x01}
\end{tblr}

The data structure of the Node Property when 0x00 (Value-Limited Purse Service) is specified as the target is shown below. (Table~\ref{table:felica-limited-purse-service-property-structure})

\begin{tblr}[
        long,
        caption={Data structure of the Node Property of a Value-Limited Purse Service},
        label={table:felica-limited-purse-service-property-structure},
    ]{
        colspec={|r|r|l|X|},
        rowhead=1,
        row{1}={c,m},
        vlines,
        hlines,
    }
    Offset & Length & Item                             & Value or remarks                                                  \\
    0x00   & 0x01   & Value-Limited Purse Service flag & 0x01: enabled, 0x00: disabled                                     \\
    0x01   & 0x04   & Upper Limit                      & Little endian, two's complement. All bytes are 0xFF when disabled \\
    0x05   & 0x04   & Lower Limit                      & Little endian, two's complement. All bytes are 0xFF when disabled \\
    0x09   & 0x01   & Generation Number                & 0xFF when disabled
\end{tblr}

When 0x01 (Communication-with-MAC-enabled Service) is specified as the target, the Node Property is a single byte holding the Communication with MAC flag. A value of 0x01 means enabled and 0x00 means disabled.

\section{Request Service v2}

The Request Service v2 command obtains the Encryption Identifier and Key Version information.

\subsection{Data Structure of the Command Packet}

The data structure of the Request Service v2 command packet is shown below. (Table~\ref{table:felica-request-service-v2-command-packet-structure})

\begin{tblr}[
        long,
        caption={Data structure of the Request Service v2 command packet},
        label={table:felica-request-service-v2-command-packet-structure},
    ]{
        colspec={|r|r|l|X|},
        rowhead=1,
        row{1}={c,m},
        vlines,
        hlines,
    }
    Offset & Length & Item               & Value or remarks          \\
    0x00   & 0x01   & Command Code       & 0x32                      \\
    0x01   & 0x08   & IDm                & -                         \\
    0x09   & 0x01   & Number of Node (n) & 0x01 $\leq$ n $\leq$ 0x20 \\
    0x0A   & 2n     & Node Code List     & Little endian
\end{tblr}

\subsection{Data Structure of the Response Packet}

The data structure of the Request Service v2 response packet is shown below. (Table~\ref{table:felica-request-service-v2-response-packet-structure})

\begin{tblr}[
        long,
        caption={Data structure of the Request Service v2 response packet},
        label={table:felica-request-service-v2-response-packet-structure},
    ]{
        colspec={|r|r|l|X|},
        rowhead=1,
        row{1}={c,m},
        vlines,
        hlines,
    }
    Offset & Length   & Item                  & Value or remarks                                                   \\
    0x00   & 0x01     & Response Code         & 0x33                                                               \\
    0x01   & 0x08     & IDm                   & -                                                                  \\
    0x09   & 0x01     & Status Flag1          & 0x00: normal                                                       \\
    0x0A   & 0x01     & Status Flag2          & 0x00: normal                                                       \\
    0x0B   & 0x01     & Encryption Identifier & Returned only when Status Flag1 is 0x00                            \\
    0x0C   & 0x01     & Number of Node (n)    & Returned only when Status Flag1 is 0x00. 0x01 $\leq$ n $\leq$ 0x20 \\
    0x0D   & 2n or 4n & Node Key Version List & {Returned only when Status Flag1 is 0x00.                          \\ 4n bytes when the Encryption Identifier is 0x41 or 0x43 (AES first, DES second). \\ 2n bytes otherwise}
\end{tblr}

\section{Get System Status}

The Get System Status command obtains the configuration state of each System.

\subsection{Data Structure of the Command Packet}

The data structure of the Get System Status command packet is shown below. (Table~\ref{table:felica-get-system-status-command-packet-structure})

\begin{tblr}[
        long,
        caption={Data structure of the Get System Status command packet},
        label={table:felica-get-system-status-command-packet-structure},
    ]{
        colspec={|r|r|l|X|},
        rowhead=1,
        row{1}={c,m},
        vlines,
        hlines,
    }
    Offset & Length & Item         & Value or remarks \\
    0x00   & 0x01   & Command Code & 0x38             \\
    0x01   & 0x08   & IDm          & -                \\
    0x09   & 0x02   & Reserved     & 0x0000
\end{tblr}

\subsection{Data Structure of the Response Packet}

The data structure of the Get System Status response packet is shown below. (Table~\ref{table:felica-get-system-status-response-packet-structure})

\begin{tblr}[
        long,
        caption={Data structure of the Get System Status response packet},
        label={table:felica-get-system-status-response-packet-structure},
    ]{
        colspec={|r|r|l|X|},
        rowhead=1,
        row{1}={c,m},
        vlines,
        hlines,
    }
    Offset & Length & Item            & Value or remarks \\
    0x00   & 0x01   & Response Code   & 0x39             \\
    0x01   & 0x08   & IDm             & -                \\
    0x09   & 0x01   & Status Flag1    & 0x00: normal     \\
    0x0A   & 0x01   & Status Flag2    & 0x00: normal     \\
    0x0B   & 0x01   & Flag            & Details unknown  \\
    0x0C   & 0x01   & Data length (n) & -                \\
    0x0D   & n      & Data            & Details unknown
\end{tblr}

The meaning of the flag and of the data is unknown.

\section{Request Specification Version}

The Request Specification Version command obtains the version of the card OS.

\subsection{Data Structure of the Command Packet}

The data structure of the Request Specification Version command packet is shown below. (Table~\ref{table:felica-request-specification-version-command-packet-structure})

\begin{tblr}[
        long,
        caption={Data structure of the Request Specification Version command packet},
        label={table:felica-request-specification-version-command-packet-structure},
    ]{
        colspec={|r|r|l|X|},
        rowhead=1,
        row{1}={c,m},
        vlines,
        hlines,
    }
    Offset & Length & Item         & Value or remarks \\
    0x00   & 0x01   & Command Code & 0x3C             \\
    0x01   & 0x08   & IDm          & -                \\
    0x09   & 0x02   & Reserved     & 0x0000
\end{tblr}

\subsection{Data Structure of the Response Packet}

The data structure of the Request Specification Version response packet is shown below. (Table~\ref{table:felica-request-specification-version-response-packet-structure})

\begin{tblr}[
        long,
        caption={Data structure of the Request Specification Version response packet},
        label={table:felica-request-specification-version-response-packet-structure},
    ]{
        colspec={|r|r|l|X|},
        rowhead=1,
        row{1}={c,m},
        vlines,
        hlines,
    }
    Offset & Length & Item                 & Value or remarks                                       \\
    0x00   & 0x01   & Response Code        & 0x3D                                                   \\
    0x01   & 0x08   & IDm                  & -                                                      \\
    0x09   & 0x01   & Status Flag1         & 0x00: normal                                           \\
    0x0A   & 0x01   & Status Flag2         & 0x00: normal                                           \\
    0x0B   & 0x01   & Format Version       & Fixed at 0x00. Returned only when Status Flag1 is 0x00 \\
    0x0C   & 0x02   & Basic Version        & Returned only when Status Flag1 is 0x00. Little endian \\
    0x0E   & 0x01   & Number of Option (m) & Returned only when Status Flag1 is 0x00                \\
    0x0F   & 2m     & Option Version List  & Returned only when Status Flag1 is 0x00. Little endian
\end{tblr}

The assignment of each element of the Option Version List is shown below. (Table~\ref{table:felica-option-version-list})

\begin{tblr}[
        long,
        caption={Assignment of the Option Version List},
        label={table:felica-option-version-list},
    ]{
        colspec={|l|X|},
        rowhead=1,
        row{1}={c,m},
        vlines,
        hlines,
    }
    Position & Option                                                       \\
    D0-D1    & DES option                                                   \\
    D2-D3    & 0x00, 0x00 (mobile FeliCa may set a version of its own here) \\
    D4-D5    & Extended Overlap option                                      \\
    D6-D7    & Value-Limited Purse Service option                           \\
    D8-D9    & Communication with MAC option                                \\
    D10-D11  & Random ID option
\end{tblr}

The Basic Version and each option version are two-byte values. The upper four bits are fixed at 0b1000 and the remaining twelve bits give the version value in BCD. For example, version 5.0.0 is 0x500. Which options are present and what their version values are differ from product to product.

\section{Reset Mode}

The Reset Mode command resets the Mode to Mode0.

\subsection{Data Structure of the Command Packet}

The data structure of the Reset Mode command packet is shown below. (Table~\ref{table:felica-reset-mode-command-packet-structure})

\begin{tblr}[
        long,
        caption={Data structure of the Reset Mode command packet},
        label={table:felica-reset-mode-command-packet-structure},
    ]{
        colspec={|r|r|l|X|},
        rowhead=1,
        row{1}={c,m},
        vlines,
        hlines,
    }
    Offset & Length & Item         & Value or remarks \\
    0x00   & 0x01   & Command Code & 0x3E             \\
    0x01   & 0x08   & IDm          & -                \\
    0x09   & 0x02   & Reserved     & 0x0000
\end{tblr}

\subsection{Data Structure of the Response Packet}

The data structure of the Reset Mode response packet is shown below. (Table~\ref{table:felica-reset-mode-response-packet-structure})

\begin{tblr}[
        long,
        caption={Data structure of the Reset Mode response packet},
        label={table:felica-reset-mode-response-packet-structure},
    ]{
        colspec={|r|r|l|X|},
        rowhead=1,
        row{1}={c,m},
        vlines,
        hlines,
    }
    Offset & Length & Item          & Value or remarks \\
    0x00   & 0x01   & Response Code & 0x3F             \\
    0x01   & 0x08   & IDm           & -                \\
    0x09   & 0x01   & Status Flag1  & 0x00: normal     \\
    0x0A   & 0x01   & Status Flag2  & 0x00: normal
\end{tblr}

\section{Authentication1 v2}

The Authentication1 v2 command starts the card-side authentication challenge under AES.

\subsection{Data Structure of the Command Packet}

For the data structure of the Authentication1 v2 command packet, see Table~\ref{table:felica-auth1v2-command-packet-structure}.

\subsection{Data Structure of the Response Packet}

For the data structure of the Authentication1 v2 response packet, see Table~\ref{table:felica-auth1v2-response-packet-structure}.

\section{Authentication2 v2}

The Authentication2 v2 command completes AES mutual authentication.

\subsection{Data Structure of the Command Packet}

For the data structure of the Authentication2 v2 command packet, see Table~\ref{table:felica-auth2v2-command-packet-structure}.

\subsection{Data Structure of the Response Packet}

For the data structure of the Authentication2 v2 response packet, see Table~\ref{table:felica-auth2v2-response-packet-structure}.

\section{Read v2}

The Read v2 command reads Block Data in an AES session established after mutual authentication.

\subsection{Data Structure of the Command Packet}

For the structure of the internal payload of the Read v2 command before encryption, see Table~\ref{table:felica-secure-read-common-command-payload-structure}.

\subsection{Data Structure of the Response Packet}

For the structure of the internal payload of the Read v2 response before encryption, see Table~\ref{table:felica-secure-read-common-response-payload-structure}.

\section{Write v2}

The Write v2 command writes Block Data in an AES session established after mutual authentication.

\subsection{Data Structure of the Command Packet}

For the structure of the internal payload of the Write v2 command before encryption, see Table~\ref{table:felica-secure-write-common-command-payload-structure}.
Key change (Access Mode 0b100) is not valid for Write v2.

\subsection{Data Structure of the Response Packet}

For the structure of the internal payload of the Write v2 response before encryption, see Table~\ref{table:felica-secure-write-common-response-payload-structure}.

\section{Register Issue ID}

The Register Issue ID command is an issuance secure command that registers information related to the issue ID.

\subsection{Data Structure of the Command Packet}

For the structure of the internal payload of the Register Issue ID command before encryption, see Table~\ref{table:felica-secure-register-issue-id-command-payload-structure}.

\subsection{Data Structure of the Response Packet}

In addition to Status Flag1 and Status Flag2, the internal payload of the Register Issue ID response before encryption returns the number of remaining Blocks (two bytes) on success only.

\section{Register Area}

The Register Area command is an issuance secure command that registers an Area.

\subsection{Data Structure of the Command Packet}

For the structure of the internal payload of the Register Area command before encryption, see Table~\ref{table:felica-secure-register-area-command-payload-structure}.

\subsection{Data Structure of the Response Packet}

The internal payload of the Register Area response before encryption is the two bytes of Status Flag1 and Status Flag2.

\section{Register Service}

The Register Service command is an issuance secure command that registers a Service.

\subsection{Data Structure of the Command Packet}

For the structure of the internal payload of the Register Service command before encryption, see Table~\ref{table:felica-secure-register-service-command-payload-structure}.

\subsection{Data Structure of the Response Packet}

In addition to Status Flag1 and Status Flag2, the internal payload of the Register Service response before encryption returns the number of remaining Blocks (two bytes) on success only.

\section{Change System Block}

The Change System Block command is an issuance secure command that commits the results of issuance commands.

\subsection{Data Structure of the Command Packet}

For the structure of the internal payload of the Change System Block command before encryption, see Table~\ref{table:felica-secure-change-system-block-command-payload-structure}.

\subsection{Data Structure of the Response Packet}

The internal payload of the Change System Block response before encryption is the two bytes of Status Flag1 and Status Flag2.

\section{Status Flag}

Status Flag is the value that represents the result of processing a command. It consists of Status Flag1 (1 byte) and Status Flag2 (1 byte).

\subsection{Status Flag1}

Status Flag1 indicates whether the command succeeded, and the Block position or Service position at which an error occurred. (Table~\ref{table:felica-status-flag-1})

\begin{tblr}[
        long,
        caption={Values and meanings of Status Flag1},
        label={table:felica-status-flag-1},
    ]{
        colspec={|r|X|},
        rowhead=1,
        row{1}={c,m},
        vlines,
        hlines,
    }
    Value & Meaning                                                                                                 \\
    0x00  & Normal completion                                                                                       \\
    0xFF  & An error in a command whose command packet contains no list, or an error that does not depend on a list \\
    Other & The position in the list at which the error occurred
\end{tblr}

There are two ways in which the position of the error is represented, depending on the product, and the value of Status Flag1 alone does not tell which one is in use. (Table~\ref{table:felica-status-flag-1-formats}) For example, when an error occurs at the Node specified tenth in the Block List, the ordinal form returns 0x0A whereas the bitmap form returns 0x02.

\begin{tblr}[
        long,
        caption={Representations of the error position in Status Flag1},
        label={table:felica-status-flag-1-formats},
    ]{
        colspec={|l|X|},
        rowhead=1,
        row{1}={c,m},
        vlines,
        hlines,
    }
    Form         & Content                                                                                                                                                                                                                                                                                     \\
    Ordinal form & The position (starting from 1) in the Block List or the Service Code List at which the error occurred is set as is                                                                                                                                                                          \\
    Bitmap form  & Each bit corresponds to a position in the list, where 1 means an error and 0 means no error. Bit 0 corresponds to the first or the ninth entry, bit 1 to the second or the tenth, and so on up to bit 6, which corresponds to the seventh or the fifteenth; bit 7 corresponds to the eighth
\end{tblr}

\subsection{Status Flag2}

Status Flag2 indicates the detail of the error. (Table~\ref{table:felica-status-flag-2}) There are also card-specific values, which are not common to all products. (Table~\ref{table:felica-status-flag-2-ex})

\begin{tblr}[
        long,
        caption={Values and meanings of Status Flag2},
        label={table:felica-status-flag-2},
    ]{
        colspec={|r|X|},
        rowhead=1,
        row{1}={c,m},
        vlines,
        hlines,
    }
    Value & Meaning                                                                                                                                            \\
    0x00  & Normal completion                                                                                                                                  \\
    0x01  & On a purse decrement the result of the calculation would be less than zero, or on a cashback the result of the calculation would exceed four bytes \\
    0x02  & On a purse cashback the specified data exceeds the value of the cashback data                                                                      \\
    0x03  & On a write to a Value-Limited Purse Service the purse data does not fall between the Upper Limit and the Lower Limit                               \\
    0x70  & Memory error (fatal error)                                                                                                                         \\
    0x71  & The number of memory rewrites has exceeded the limit (a warning; the write is still performed)
\end{tblr}

Because 0x71 is a warning, the write itself is performed. The limit on the number of rewrites differs from product to product, and in this case Status Flag1 is 0x00 on some products and 0xFF on others.

\begin{tblr}[
        long,
        caption={Values and meanings of Status Flag2 (card-specific)},
        label={table:felica-status-flag-2-ex},
    ]{
        colspec={|r|X|},
        rowhead=1,
        row{1}={c,m},
        vlines,
        hlines,
    }
    Value & Meaning                                                                                                                                                                     \\
    0xA1  & The number of Services or the number of Nodes specified in the command is outside the defined range                                                                         \\
    0xA2  & The number of Blocks specified in the command is outside the range defined for the product                                                                                  \\
    0xA3  & The Service Code List Order specified in a Block List Element is outside the range of the number of Services specified in the command or specified at mutual authentication \\
    0xA4  & The Area Attribute of the Area Code or the Service Attribute of the Service Code specified in the command is wrong                                                          \\
    0xA5  & The Area or the Service specified in the command cannot be accessed, or a parameter specified in the command does not satisfy the success conditions                        \\
    0xA6  & The access target identified by the Service Code List Order specified in a Block List Element, or the Node specified in the Node Code List, does not exist                  \\
    0xA7  & The Access Mode specified in a Block List Element is wrong                                                                                                                  \\
    0xA8  & The Block Number specified in a Block List Element exceeds the number of Blocks allocated to the Service                                                                    \\
    0xA9  & A write failed during an issuance command                                                                                                                                   \\
    0xAA  & Key change failed                                                                                                                                                           \\
    0xAB  & The package parity or the package MAC is invalid during an issuance command                                                                                                 \\
    0xAC  & A parameter is invalid during an issuance command                                                                                                                           \\
    0xAD  & The Service to be registered already exists                                                                                                                                 \\
    0xAE  & The System Code is invalid during an issuance command                                                                                                                       \\
    0xAF  & The number of Blocks written simultaneously to the cyclic Service specified in the command exceeds the number of Blocks allocated to the Service                            \\
    0xC0  & The package identifier is invalid during an issuance command                                                                                                                \\
    0xC1  & The parameters inside and outside the package do not match during an issuance command                                                                                       \\
    0xC2  & The issuance command has been disabled                                                                                                                                      \\
    0xC3  & The attribute of the Node specified in the command is wrong
\end{tblr}

\chapter{Mutual Authentication and Secure Messaging Algorithms}

\section{Authentication1 / Authentication2 (DES)}

\subsection{Data Structure of the Command Packets}

The data structure of the Authentication1 command packet is shown below. (Table~\ref{table:felica-auth1-command-packet-structure})

\begin{tblr}[
        long,
        caption={Data structure of the Authentication1 command packet},
        label={table:felica-auth1-command-packet-structure},
    ]{
        colspec={|r|r|l|X|},
        rowhead=1,
        row{1}={c,m},
        vlines,
        hlines,
    }
    Offset & Length & Item                  & Value or remarks \\
    0x00   & 0x01   & Command Code          & 0x10             \\
    0x01   & 0x08   & IDm                   & -                \\
    0x09   & 0x01   & Number of Areas (a)   & -                \\
    0x0A   & 2a     & Area Code list        & Little endian    \\
    -      & 0x01   & Number of Service (s) & -                \\
    -      & 2s     & Service Code List     & Little endian    \\
    -      & 0x08   & Challenge 1A          & 8 bytes
\end{tblr}

The data structure of the Authentication2 command packet is shown below. (Table~\ref{table:felica-auth2-command-packet-structure})

\begin{tblr}[
        long,
        caption={Data structure of the Authentication2 command packet},
        label={table:felica-auth2-command-packet-structure},
    ]{
        colspec={|r|r|l|X|},
        rowhead=1,
        row{1}={c,m},
        vlines,
        hlines,
    }
    Offset & Length & Item         & Value or remarks \\
    0x00   & 0x01   & Command Code & 0x12             \\
    0x01   & 0x08   & IDm          & -                \\
    0x09   & 0x08   & Challenge 2B & 8 bytes
\end{tblr}

\subsection{Data Structure of the Response Packets}

The data structure of the Authentication1 response packet is shown below. (Table~\ref{table:felica-auth1-response-packet-structure})

\begin{tblr}[
        long,
        caption={Data structure of the Authentication1 response packet},
        label={table:felica-auth1-response-packet-structure},
    ]{
        colspec={|r|r|l|X|},
        rowhead=1,
        row{1}={c,m},
        vlines,
        hlines,
    }
    Offset & Length & Item          & Value or remarks \\
    0x00   & 0x01   & Response Code & 0x11             \\
    0x01   & 0x08   & IDm           & -                \\
    0x09   & 0x08   & Challenge 1B  & 8 bytes          \\
    0x11   & 0x08   & Challenge 2A  & 8 bytes
\end{tblr}

The data structure of the Authentication2 response packet is shown below. (Table~\ref{table:felica-auth2-response-packet-structure}) This response contains no IDm, and everything after the Response Code is encrypted.

\begin{tblr}[
        long,
        caption={Data structure of the Authentication2 response packet},
        label={table:felica-auth2-response-packet-structure},
    ]{
        colspec={|r|r|l|X|},
        rowhead=1,
        row{1}={c,m},
        vlines,
        hlines,
    }
    Offset & Length & Item              & Value or remarks                                                                                                                                                                \\
    0x00   & 0x01   & Response Code     & 0x13                                                                                                                                                                            \\
    0x01   & 0x20   & Encrypted payload & Encrypted with DES-CBC (IV = 0) keyed with the random number $R_2$. For the structure after decryption, see Table~\ref{table:felica-auth2-response-decrypted-payload-structure}
\end{tblr}

The data structure of the payload after decryption is shown below. (Table~\ref{table:felica-auth2-response-decrypted-payload-structure})


\begin{tblr}[
        long,
        caption={Data structure of the decrypted payload of the Authentication2 response},
        label={table:felica-auth2-response-decrypted-payload-structure},
    ]{
        colspec={|r|r|l|X|},
        rowhead=1,
        row{1}={c,m},
        vlines,
        hlines,
    }
    Offset & Length & Item                    & Value or remarks                                      \\
    0x00   & 0x02   & Transaction number (TN) & Little endian                                         \\
    0x02   & 0x06   & Transaction ID (TID)    & Matches the last six bytes of the random number $R_1$ \\
    0x08   & 0x08   & Issue ID (IDi)          & -                                                     \\
    0x10   & 0x08   & Issue Parameter (PMi)   & -                                                     \\
    0x18   & 0x08   & MAC                     & Computed over the preceding 24 bytes
\end{tblr}

\subsection{DES Mutual Authentication Algorithm}

The principal parameters used in DES mutual authentication are derived as follows.

\begin{itemize}
    \item Let the System Key be $K_{\mathrm{sys}}$, the sequence of Area Keys be $K_{\mathrm{Area},i}$ and the sequence of Service Keys be $K_{\mathrm{srv},j}$. Then
          \[
              G_0 = K_{\mathrm{sys}},\quad
              G_i = \mathrm{DES}_{K_{\mathrm{Area},i}}(G_{i-1})
          \]
          \[
              K_{\mathrm{group}} = G_m,\quad
              U_0 = K_{\mathrm{group}},\quad
              U_j = \mathrm{DES}_{K_{\mathrm{srv},j}}(U_{j-1}),\quad
              K_{\mathrm{user}} = U_n
          \]
          give the group Service Key $K_{\mathrm{group}}$ and the user Service Key $K_{\mathrm{user}}$.
    \item Treating the IDm as an eight-byte vector, compute
          \[
              L = K_{\mathrm{group}} \oplus IDm
          \]
          \[
              \alpha = \mathrm{DES}_{L}(K_{\mathrm{user}}),\quad
              \beta  = \mathrm{DES}_{\alpha}(L)
          \]
    \item Random numbers are exchanged with two-key 3DES, in the key order $K_1$, $K_2$, $K_1$. Letting $R_1$ and $R_2$ be the random numbers,
          \[
              C_{1A} = \mathrm{3DES}_{\alpha,L}(R_1),\quad
              C_{1B} = \mathrm{3DES}_{L,\beta}(R_1)
          \]
          \[
              C_{2A} = \mathrm{3DES}_{L,\beta}(R_2),\quad
              C_{2B} = \mathrm{3DES}_{\alpha,L}(R_2)
          \]
    \item The plaintext of Authentication2 is
          \[
              \mathrm{TN}(2)\ \|\ \mathrm{TID}(6)\ \|\ \mathrm{IDi}(8)\ \|\ \mathrm{PMi}(8)
          \]
          Here the TID is the last six bytes of $R_1$. The data used as the Issue ID information is the last sixteen bytes, $\mathrm{IDi}(8)\ \|\ \mathrm{PMi}(8)$.
\end{itemize}

\section[Authentication1 v2 / Authentication2 v2 (AES)]{Authentication1 v2 /\\ Authentication2 v2 (AES)}

\subsection{Data Structure of the Command Packets}

The data structure of the Authentication1 v2 command packet is shown below. (Table~\ref{table:felica-auth1v2-command-packet-structure})

\begin{tblr}[
        long,
        caption={Data structure of the Authentication1 v2 command packet},
        label={table:felica-auth1v2-command-packet-structure},
    ]{
        colspec={|r|r|l|X|},
        rowhead=1,
        row{1}={c,m},
        vlines,
        hlines,
    }
    Offset & Length & Item                & Value or remarks \\
    0x00   & 0x01   & Command Code        & 0x40             \\
    0x01   & 0x08   & IDm                 & -                \\
    0x09   & 0x01   & Operation Parameter & -                \\
    0x0A   & 0x01   & Number of Node (n)  & -                \\
    0x0B   & 2n     & Node Code List      & Little endian    \\
    -      & 0x10   & Challenge 1A        & 16 bytes
\end{tblr}

The data structure of the Authentication2 v2 command packet is shown below. (Table~\ref{table:felica-auth2v2-command-packet-structure})

\begin{tblr}[
        long,
        caption={Data structure of the Authentication2 v2 command packet},
        label={table:felica-auth2v2-command-packet-structure},
    ]{
        colspec={|r|r|l|X|},
        rowhead=1,
        row{1}={c,m},
        vlines,
        hlines,
    }
    Offset & Length & Item         & Value or remarks \\
    0x00   & 0x01   & Command Code & 0x42             \\
    0x01   & 0x08   & IDm          & -                \\
    0x09   & 0x10   & Challenge 2B & 16 bytes
\end{tblr}

\subsection{Data Structure of the Response Packets}

The data structure of the Authentication1 v2 response packet is shown below. (Table~\ref{table:felica-auth1v2-response-packet-structure})

\begin{tblr}[
        long,
        caption={Data structure of the Authentication1 v2 response packet},
        label={table:felica-auth1v2-response-packet-structure},
    ]{
        colspec={|r|r|l|X|},
        rowhead=1,
        row{1}={c,m},
        vlines,
        hlines,
    }
    Offset & Length & Item          & Value or remarks \\
    0x00   & 0x01   & Response Code & 0x41             \\
    0x01   & 0x08   & IDm           & -                \\
    0x09   & 0x10   & Challenge 1B  & 16 bytes         \\
    0x19   & 0x10   & Challenge 2A  & 16 bytes         \\
    0x29   & 0x04   & Challenge 3C  & 4 bytes
\end{tblr}

The data structure of the Authentication2 v2 response packet is shown below. (Table~\ref{table:felica-auth2v2-response-packet-structure}) This response contains no IDm either. Only the transaction number is in cleartext; the payload and the MAC are encrypted.

\begin{tblr}[
        long,
        caption={Data structure of the Authentication2 v2 response packet},
        label={table:felica-auth2v2-response-packet-structure},
    ]{
        colspec={|r|r|l|X|},
        rowhead=1,
        row{1}={c,m},
        vlines,
        hlines,
    }
    Offset & Length & Item                    & Value or remarks                                                                                                                          \\
    0x00   & 0x01   & Response Code           & 0x43                                                                                                                                      \\
    0x01   & 0x02   & Transaction number (TN) & Little endian. Not encrypted                                                                                                              \\
    0x03   & 0x10   & Encrypted payload       & Encrypted with AES-128-OFB. For the structure after decryption, see Table~\ref{table:felica-auth2v2-response-decrypted-payload-structure} \\
    0x13   & 0x08   & Encrypted MAC           & Encrypted with the later part of the same OFB stream
\end{tblr}

The data structure of the payload after decryption is shown below. (Table~\ref{table:felica-auth2v2-response-decrypted-payload-structure})


\begin{tblr}[
        long,
        caption={Data structure of the decrypted payload of the Authentication2 v2 response},
        label={table:felica-auth2v2-response-decrypted-payload-structure},
    ]{
        colspec={|r|r|l|X|},
        rowhead=1,
        row{1}={c,m},
        vlines,
        hlines,
    }
    Offset & Length & Item                  & Value or remarks \\
    0x00   & 0x08   & Issue ID (IDi)        & -                \\
    0x08   & 0x08   & Issue Parameter (PMi) & -
\end{tblr}

\subsection{AES Mutual Authentication Algorithm}

The principal parameters used in AES mutual authentication are derived as follows.

\begin{itemize}
    \item Let the sixteen-byte constant be $I_0 = (01\ 01\ 00\cdots00\ 80)$ and the sequence of Node Keys be $K_{\mathrm{Node},j}$. Then
          \[
              S_0 = I_0,\quad
              S_j = \mathrm{AES}_{K_{\mathrm{Node},j}}(S_{j-1}),\quad
              K_{\mathrm{group}} = S_n
          \]
          gives the group key $K_{\mathrm{group}}$.
    \item With the group key $K_{\mathrm{group}}$ and the individualisation code $K_{\mathrm{ind}}$, compute
          \[
              H = K_{\mathrm{group}} \oplus K_{\mathrm{ind}}
          \]
    \item Define the sixteen-byte context Block $B(\mathrm{prefix},IDm)$ as
          \[
              B[0..1]=\mathrm{prefix},\
              B[2..5]=0,\
              B[6..13]=IDm,\
              B[14..15]=0x01\ 0x00
          \]
          and obtain
          \[
              \alpha = \mathrm{AES}_{H}(B([0x01,0x02],IDm))
          \]
          \[
              \beta  = \mathrm{AES}_{H}(B([0x02,0x02],IDm))
          \]
    \item From challenge 3C (four bytes), form
          \[
              \beta' = \beta \oplus (C_{3C}\ \|\ 0x00^{12})
          \]
          and generate and verify the authentication values with
          \[
              C_{1A}=\mathrm{AES}_{\alpha}(R_1),\quad
              C_{1B}=\mathrm{AES}_{\beta'}(R_1)
          \]
          \[
              C_{2A}=\mathrm{AES}_{\beta'}(R_2),\quad
              C_{2B}=\mathrm{AES}_{\alpha}(R_2)
          \]
    \item The TID used in secure messaging is derived from $R_1$ as
          \[
              \mathrm{TID}=R_1[2..7]
          \]
          (six bytes, from the third to the eighth byte, counting the first byte as 0).
    \item The decrypted payload of the Authentication2 v2 response is
          \[
              \mathrm{IDi}(8)\ \|\ \mathrm{PMi}(8)
          \]
    \item After Authentication2 v2 succeeds, the session keys are derived from $R_2$.
          \[
              K_{\mathrm{enc}}=\mathrm{AES}_{R_2}(01\ 03\ 00\cdots00\ 01\ 00)
          \]
          \[
              K_{\mathrm{mac}}=\mathrm{AES}_{R_2}(02\ 03\ 00\cdots00\ 01\ 00)
          \]
\end{itemize}

\section{Secure Read/Write Command Format}

Read, Write, Read v2 and Write v2 carry internal payloads of identical form inside the encrypted secure frame. The only difference is the encryption scheme.

\begin{tblr}[
        long,
        caption={Differences between the four secure Read/Write commands},
        label={table:felica-secure-read-write-command-diff},
    ]{
        colspec={|l|l|X|},
        rowhead=1,
        row{1}={c,m},
        vlines,
        hlines,
    }
    Command  & Scheme & Payload structure                                                                                                                                                               \\
    Read     & DES    & Common to the Read family (Table~\ref{table:felica-secure-read-common-command-payload-structure} and Table~\ref{table:felica-secure-read-common-response-payload-structure})    \\
    Read v2  & AES    & Common to the Read family (Table~\ref{table:felica-secure-read-common-command-payload-structure} and Table~\ref{table:felica-secure-read-common-response-payload-structure})    \\
    Write    & DES    & Common to the Write family (Table~\ref{table:felica-secure-write-common-command-payload-structure} and Table~\ref{table:felica-secure-write-common-response-payload-structure}) \\
    Write v2 & AES    & Common to the Write family (Table~\ref{table:felica-secure-write-common-command-payload-structure} and Table~\ref{table:felica-secure-write-common-response-payload-structure})
\end{tblr}

\begin{tblr}[
        long,
        caption={Data structure of the internal payload of the Read family commands (Read/Read v2)},
        label={table:felica-secure-read-common-command-payload-structure},
    ]{
        colspec={|r|r|l|X|},
        rowhead=1,
        row{1}={c,m},
        vlines,
        hlines,
    }
    Offset & Length   & Item                & Value or remarks                                                                         \\
    0x00   & 0x01     & Number of Block (n) & 0x01 $\leq$ n $\leq$ the maximum number of simultaneously readable Blocks of the product \\
    0x01   & 2n to 3n & Block List          & Two-byte or three-byte form of the Block List Element
\end{tblr}

\begin{tblr}[
        long,
        caption={Data structure of the internal payload of the Read family responses (Read/Read v2)},
        label={table:felica-secure-read-common-response-payload-structure},
    ]{
        colspec={|r|r|l|X|},
        rowhead=1,
        row{1}={c,m},
        vlines,
        hlines,
    }
    Offset & Length & Item                & Value or remarks                        \\
    0x00   & 0x01   & Status Flag1        & 0x00: normal                            \\
    0x01   & 0x01   & Status Flag2        & 0x00: normal                            \\
    0x02   & 0x01   & Number of Block (n) & Returned only when Status Flag1 is 0x00 \\
    0x03   & 16n    & Block Data          & Returned only when Status Flag1 is 0x00
\end{tblr}

\begin{tblr}[
        long,
        caption={Data structure of the internal payload of the Write family commands (Write/Write v2)},
        label={table:felica-secure-write-common-command-payload-structure},
    ]{
        colspec={|r|r|l|X|},
        rowhead=1,
        row{1}={c,m},
        vlines,
        hlines,
    }
    Offset   & Length   & Item                & Value or remarks                                                                         \\
    0x00     & 0x01     & Number of Block (n) & 0x01 $\leq$ n $\leq$ the maximum number of simultaneously writable Blocks of the product \\
    0x01     & 2n to 3n & Block List          & Two-byte or three-byte form of the Block List Element                                    \\
    Variable & 16n      & Block Data          & 16 bytes $\times$ n
\end{tblr}

\begin{tblr}[
        long,
        caption={Data structure of the internal payload of the Write family responses (Write/Write v2)},
        label={table:felica-secure-write-common-response-payload-structure},
    ]{
        colspec={|r|r|l|X|},
        rowhead=1,
        row{1}={c,m},
        vlines,
        hlines,
    }
    Offset & Length & Item         & Value or remarks \\
    0x00   & 0x01   & Status Flag1 & 0x00: normal     \\
    0x01   & 0x01   & Status Flag2 & 0x00: normal
\end{tblr}

\subsection{Generating and Using the DES Key Change Package}

To perform a DES key change with the Write command, set the Access Mode of the Block List Element to 0b100 (key change) and store the sixteen-byte key change package in the Block Data.

\begin{tblr}[
        long,
        caption={Data structure of the key change package (16 bytes)},
        label={table:felica-secure-key-change-package-structure},
    ]{
        colspec={|r|r|l|X|},
        rowhead=1,
        row{1}={c,m},
        vlines,
        hlines,
    }
    Offset & Length & Item       & Value or remarks                                        \\
    0x00   & 0x08   & Parameter1 & Cipher Block containing the new Key Version information \\
    0x08   & 0x08   & Parameter2 & Cipher Block containing the new key itself
\end{tblr}

Let the parent key be $K_{\mathrm{parent}}$, the old key $K_{\mathrm{old}}$, the new key $K_{\mathrm{new}}$ and the new Key Version $v$. First build the eight-byte version Block $V$ as
\[
    V = 00\ 00\ 00\ 00\ 00\ 00\ v_{\mathrm{LE},0}\ v_{\mathrm{LE},1}
\]
then compute the following.
\[
    P_1 = \mathrm{DES}_{K_{\mathrm{parent}}}\Bigl(
    \mathrm{DES}_{K_{\mathrm{old}}}\bigl(
    \mathrm{DES}_{K_{\mathrm{new}}}(V)\bigr)\Bigr)
\]
\[
    P_2 = \mathrm{DES}_{K_{\mathrm{parent}}}\Bigl(
    \mathrm{DES}_{K_{\mathrm{old}}}(K_{\mathrm{new}})\Bigr)
\]
\[
    \mathrm{KeyChangePackage} = P_1\ \|\ P_2
\]

In use, the following pair is passed to the Write command for each key to be changed.
\begin{itemize}
    \item Block List Element: Access Mode = 0b100, Block Number / Key Version = $v$, Service Code List Order = the target Service
    \item Block Data (16 bytes): the $\mathrm{KeyChangePackage}$ above
\end{itemize}

When changing several keys at once, arrange the Block List and the Block Data so that their orders match. On failure, Status Flag2 may be, for example, 0xAA (key change failed) or 0xAB (invalid package MAC).

\subsection{DES Secure Messaging}

In a DES session, the payload before encryption is
\[
    P = \mathrm{TN}(2,\mathrm{LE})\ \|\ \mathrm{TID}(6)\ \|\ \mathrm{CommandPayload}
\]
This is padded to an eight-byte boundary by any method to give $P'$, a MAC is appended, and the result is encrypted with DES-CBC (IV = 0).

The MAC is computed as follows.
\[
    M_0 = [\mathrm{Len},\ \mathrm{Code},\ 0,\ 0,\ 0,\ 0,\ 0,\ 0]
\]
\[
    M_i = \mathrm{DES}_{B_i}(M_{i-1})
\]
Here $B_i$ is the sequence of eight-byte Blocks of $P'$ and $\mathrm{Len}=2+|P'|+8$. The final value $M_n$ is appended as the MAC (eight bytes).

The response is decrypted, its MAC is verified, the trailing eight-byte MAC is removed and then the padding is removed. The TID must also match and the TN must increase monotonically.

\subsection{AES-128 Secure Messaging}

In an AES session, the transmitted data is
\[
    \mathrm{TN}(2,\mathrm{LE})\ \|\ \mathrm{EncPayload}\ \|\ \mathrm{EncMAC}(8)
\]
The payload itself is encrypted with OFB, and the MAC is encrypted with the later part of the same OFB stream.

The initialisation vector IV (sixteen bytes) is built as follows.
\[
    IV[0]=0x01,\ IV[1]=\mathrm{FrameLen},\ IV[2]=\mathrm{Code},
\]
\[
    IV[3..4]=\mathrm{TN},\ IV[5..10]=\mathrm{TID},
\]
\[
    IV[11..13]=\mathrm{C_{3C}}[1..3],\
    IV[14..15]=0
\]

The MAC plaintext is computed with AES-CMAC, of which the first eight bytes are used.
\[
    B_0[0]=0x19,\ B_0[1..13]=IV[1..13],\ B_0[14..15]=|Payload|_{\mathrm{BE16}}
\]
\[
    MAC = \mathrm{CMAC}_{K_{\mathrm{mac}}}(B_0\ \|\ Payload)[0..7]
\]

During OFB encryption, if the payload length is not a multiple of sixteen, the keystream is consumed up to the next sixteen-byte boundary before the MAC is encrypted.

\section{DES Issuance Secure Command Format}

Register Issue ID, Register Area, Register Service and Change System Block are sent and received as DES secure commands. The structures of their internal payloads before encryption are shown below.

\begin{tblr}[
        long,
        caption={Data structure of the internal payload of the Register Issue ID command},
        label={table:felica-secure-register-issue-id-command-payload-structure},
    ]{
        colspec={|r|r|l|X|},
        rowhead=1,
        row{1}={c,m},
        vlines,
        hlines,
    }
    Offset & Length   & Item            & Value or remarks                    \\
    0x00   & 0x08     & Issue ID        & -                                   \\
    0x08   & 0x08     & Issue Parameter & -                                   \\
    0x10   & Variable & Package         & Issuance package generated with DES
\end{tblr}

\begin{tblr}[
        long,
        caption={Data structure of the internal payload of the Register Area command},
        label={table:felica-secure-register-area-command-payload-structure},
    ]{
        colspec={|r|r|l|X|},
        rowhead=1,
        row{1}={c,m},
        vlines,
        hlines,
    }
    Offset & Length   & Item      & Value or remarks                    \\
    0x00   & 0x02     & Area Code & Little endian                       \\
    0x02   & Variable & Package   & Issuance package generated with DES
\end{tblr}

\begin{tblr}[
        long,
        caption={Data structure of the internal payload of the Register Service command},
        label={table:felica-secure-register-service-command-payload-structure},
    ]{
        colspec={|r|r|l|X|},
        rowhead=1,
        row{1}={c,m},
        vlines,
        hlines,
    }
    Offset & Length   & Item         & Value or remarks                    \\
    0x00   & 0x02     & Service Code & Little endian                       \\
    0x02   & Variable & Package      & Issuance package generated with DES
\end{tblr}

\begin{tblr}[
        long,
        caption={Data structure of the internal payload of the Change System Block command},
        label={table:felica-secure-change-system-block-command-payload-structure},
    ]{
        colspec={|r|r|l|X|},
        rowhead=1,
        row{1}={c,m},
        vlines,
        hlines,
    }
    Offset & Length & Item & Value or remarks \\
    0x00   & 0x00   & None & No payload
\end{tblr}

The data structures of the internal payloads of the responses are shown below.

\begin{tblr}[
        long,
        caption={Data structure of the internal payload of the Register Issue ID response},
        label={table:felica-secure-register-issue-id-response-payload-structure},
    ]{
        colspec={|r|r|l|X|},
        rowhead=1,
        row{1}={c,m},
        vlines,
        hlines,
    }
    Offset & Length & Item                       & Value or remarks                                       \\
    0x00   & 0x01   & Status Flag1               & 0x00: normal                                           \\
    0x01   & 0x01   & Status Flag2               & 0x00: normal                                           \\
    0x02   & 0x02   & Number of remaining Blocks & Returned only when Status Flag1 is 0x00. Little endian
\end{tblr}

\begin{tblr}[
        long,
        caption={Data structure of the internal payload of the Register Area response},
        label={table:felica-secure-register-area-response-payload-structure},
    ]{
        colspec={|r|r|l|X|},
        rowhead=1,
        row{1}={c,m},
        vlines,
        hlines,
    }
    Offset & Length & Item         & Value or remarks \\
    0x00   & 0x01   & Status Flag1 & 0x00: normal     \\
    0x01   & 0x01   & Status Flag2 & 0x00: normal
\end{tblr}

\begin{tblr}[
        long,
        caption={Data structure of the internal payload of the Register Service response},
        label={table:felica-secure-register-service-response-payload-structure},
    ]{
        colspec={|r|r|l|X|},
        rowhead=1,
        row{1}={c,m},
        vlines,
        hlines,
    }
    Offset & Length & Item                       & Value or remarks                                       \\
    0x00   & 0x01   & Status Flag1               & 0x00: normal                                           \\
    0x01   & 0x01   & Status Flag2               & 0x00: normal                                           \\
    0x02   & 0x02   & Number of remaining Blocks & Returned only when Status Flag1 is 0x00. Little endian
\end{tblr}

\begin{tblr}[
        long,
        caption={Data structure of the internal payload of the Change System Block response},
        label={table:felica-secure-change-system-block-response-payload-structure},
    ]{
        colspec={|r|r|l|X|},
        rowhead=1,
        row{1}={c,m},
        vlines,
        hlines,
    }
    Offset & Length & Item         & Value or remarks \\
    0x00   & 0x01   & Status Flag1 & 0x00: normal     \\
    0x01   & 0x01   & Status Flag2 & 0x00: normal
\end{tblr}

The plaintext of an issuance package (before encryption) is sixteen bytes long. Its format for each command is as follows.

\begin{tblr}[
        long,
        caption={Data structure of the issuance package plaintext for Register Issue ID (16 bytes)},
        label={table:felica-register-issue-id-package-plain-structure},
    ]{
        colspec={|r|r|l|X|},
        rowhead=1,
        row{1}={c,m},
        vlines,
        hlines,
    }
    Offset & Length & Item               & Value or remarks \\
    0x00   & 0x02   & System Code        & Big endian       \\
    0x02   & 0x02   & Area 0 Key Version & Little endian    \\
    0x04   & 0x08   & Area 0 Key         & -                \\
    0x0C   & 0x04   & Reserved           & 0x00000000
\end{tblr}

\begin{tblr}[
        long,
        caption={Data structure of the issuance package plaintext for Register Area (16 bytes)},
        label={table:felica-register-area-package-plain-structure},
    ]{
        colspec={|r|r|l|X|},
        rowhead=1,
        row{1}={c,m},
        vlines,
        hlines,
    }
    Offset & Length & Item               & Value or remarks \\
    0x00   & 0x02   & Service Code start & Little endian    \\
    0x02   & 0x02   & Service Code end   & Little endian    \\
    0x04   & 0x02   & Size               & Little endian    \\
    0x06   & 0x02   & Key Version        & Little endian    \\
    0x08   & 0x08   & Area Key           & -
\end{tblr}

\begin{tblr}[
        long,
        caption={Data structure of the issuance package plaintext for Register Service (16 bytes)},
        label={table:felica-register-service-package-plain-structure},
    ]{
        colspec={|r|r|l|X|},
        rowhead=1,
        row{1}={c,m},
        vlines,
        hlines,
    }
    Offset & Length & Item         & Value or remarks \\
    0x00   & 0x02   & Service Code & Little endian    \\
    0x02   & 0x02   & Reserved     & 0x0000           \\
    0x04   & 0x02   & Size         & Little endian    \\
    0x06   & 0x02   & Key Version  & Little endian    \\
    0x08   & 0x08   & Service Key  & -
\end{tblr}

\subsection{Issuance Package Generation Algorithm}

Let the package key be $K_{\mathrm{pkg}}$ and the plaintext package be $P$ (where $|P|$ must be a non-zero multiple of eight).

\begin{itemize}
    \item Build the MAC generation key as
          \[
              K_{\mathrm{macpkg}} = K_{\mathrm{pkg}} \oplus 0xFF\ldots FF
          \]
    \item Compute
          \[
              C = \mathrm{DES\mbox{-}CBC}_{IV=0,K_{\mathrm{macpkg}}}(P)
          \]
          and take the last eight bytes as $MAC_{\mathrm{pkg}}$.
    \item Compute
          \[
              Package = \mathrm{DES\mbox{-}CBC}_{IV=0,K_{\mathrm{pkg}}}(P\ \|\ MAC_{\mathrm{pkg}})
          \]
          and store this ciphertext in the package field of the command.
\end{itemize}

\printbibliography[title={References}]

\end{document}
